\section{Nonlinear Elliptic Equations}\label{chap:nonlinear-elliptic-equations}

In this chapter we shall discuss nonlinear elliptic equations from three perspectives: the implicit function theorem, the calculus of variations, and nonlinear operator theory. This is the only chapter of the book in which we assume that the reader is familiar with the basic results of measure theory. In particular, we shall assume that the reader understands the following concepts and results.\begin{itemize}\item The definition of a set of measure zero and the idea of functions agreeing ``almost everywhere.\textquotesingle\textquotesingle\item The idea of Lebesgue measurable functions and the definition of the $L^p$ spaces as equivalence classes of functions that agree almost everywhere.\item The equivalence of the ``measure theoretic\textquotesingle\textquotesingle\ definition of the $L^p$ spaces and the ``completion\textquotesingle\textquotesingle\ definition used in the rest of this book.\item The idea of almost everywhere convergence of sequences of functions, the interrelationship between various types of convergence. This includes an understanding of such results as Fatou\textquotesingle s lemma and the Lebesgue dominated convergence theorem.\end{itemize}

\subsection{Perturbation Results}\label{sec:nonlinear-elliptic-perturbation-results}

Many results on differential equations say that a nonlinear equation behaves essentially like its linearization as long as one considers solutions which are small enough so that the linear terms dominate over the nonlinear ones. In Chapter~\ref{chap:introduction}, we stated the implicit function theorem from classical calculus, which provides such a result for finite-dimensional systems of equations. In this section, we shall generalize the implicit function theorem to a Banach space setting and then consider applications to elliptic PDEs.

\subsubsection{The Banach Contraction Principle and the Implicit Function Theorem}\label{subsec:banach-contraction-implicit-function}

The Banach contraction principle is one of the most used techniques for finding solutions of nonlinear equations. It consists of the following theorem.

\begin{Theorem}{Banach contraction}
\label{thm:banach-contraction-fixed-point}
Let $(X,d)$ be a complete metric space. Assume that $X$ is not empty and let $T:X\to X$ be a contraction, i.e., a mapping with the property that there exists $\theta\in[0,1)$ with the property that $d(T(x),T(y))\leq\theta d(x,y)$ for every $x,y\in X$. Then $T$ has a unique fixed point.
\end{Theorem}

\begingroup
\renewcommand{\endProof}{\par\nobreak\hfill$\blacksquare$\par\vspace{0.1cm}}
\begin{Proof}
Suppose there were two fixed points $x$ and $y$. Then
\begin{equation}\label{eq:banach-contraction-fixed-point-uniqueness}
d(x,y)=d(T(x),T(y))\leq\theta d(x,y),
\end{equation}
which is possible only if $x=y$.

Now consider any point $x_0\in X$ and define recursively $x_{n+1}=T(x_n)$. Then we have $d(x_{n+1},x_n)\leq\theta d(x_n,x_{n-1})$, from which it easily follows that the sequence $x_n$ is Cauchy. Let $x$ be the limit. Then we have
\begin{equation}\label{eq:banach-contraction-iteration-limit-fixed-point}
T(x)=T\left(\lim_{n\to\infty}x_n\right)=\lim_{n\to\infty}T(x_n)=\lim_{n\to\infty}x_{n+1}=x.
\end{equation}
This completes the proof.
\end{Proof}
\endgroup
